\section{Results}

\subsection{Deterministic validation}

All 47 tests pass, including the complete 37-test regression suite and the ten ancilla-specific tests. Python bytecode compilation and Git whitespace validation also pass. No previous target, canonicalization identity, search invariant, or certification test was modified to accommodate the new contract.

The contract tests demonstrate the intended distinction between full-unitary and isometry equality. In particular, a $Z$ acting only on a clean $\lvert0\rangle$ ancilla is accepted as contract-equivalent to identity, while a witness that leaves clean workspace populated or entangled is rejected. Borrowed workspace is restored on the entire borrowed input space, and exact mode rejects a phase difference accepted by projective mode.

\subsection{Marginal training overhead}

Table~\ref{tab:training-time} and Fig.~\ref{fig:ancilla-training} report the matched low-cost training protocol. The pre-ancilla mixed hierarchy requires \SI{2.134339}{\second}. Warm-started ancilla fine-tuning adds \SI{0.174266}{\second}, giving \SI{2.308604}{\second} in total.

\begin{table}[t]
\caption{Measured linear-policy training time.}
\label{tab:training-time}
\centering
\begin{tabular}{lr}
\toprule
Training component & Wall time (s)\\
\midrule
Existing mixed hierarchy & 2.134339\\
Ancilla-specific increment & 0.174266\\
Total ancilla-aware protocol & 2.308604\\
\midrule
Total/reference ratio & $1.08165\times$\\
Incremental fraction & $8.16\%$\\
\bottomrule
\end{tabular}
\end{table}

\begin{figure}[t]
\centering
\includegraphics[width=\columnwidth]{ancilla_training_time.png}
\caption{Matched training cost. The ancilla-aware outer and inner policies are warm-started from the existing linear models and require only a short staged fine-tuning pass.}
\label{fig:ancilla-training}
\end{figure}

The measured increase is approximately $8.2\%$, satisfying the intended marginal-overhead requirement in this qualification environment. This timing is environment-specific; the structural claim is that no new neural model, replay buffer, or joint two-policy training loop is introduced.

\subsection{Held-out clean-ancilla synthesis}

All three held-out contracts are certified. Table~\ref{tab:ancilla-results} gives the complete frozen-policy outcomes. The projective isometry error and clean-workspace leakage are zero to reported precision for every result.

\begin{table*}[t]
\caption{Frozen hierarchical synthesis under one-clean-ancilla contracts.}
\label{tab:ancilla-results}
\centering
\begin{tabular}{lrrrrrl}
\toprule
Target & Certified & Wall (s) & Allocations & Exact edges & Peak frontier & Returned witness\\
\midrule
H--T echo, 1 logical qubit
& yes & 0.001870 & 2 & 6 & 7 & $H(0);T(0)$\\
S--T product, 2 logical qubits
& yes & 0.002029 & 2 & 5 & 6 & $S(1);T(0)$\\
Mixed product, 3 logical qubits
& yes & 0.004378 & 3 & 9 & 10 & $H(2);S(0);T^\dagger(1)$\\
\bottomrule
\end{tabular}
\end{table*}

\begin{figure}[t]
\centering
\includegraphics[width=\columnwidth]{ancilla_evaluation_edges.png}
\caption{Exact continuation attempts for the three held-out contract targets.}
\label{fig:ancilla-edges}
\end{figure}

The returned circuits happen not to touch the declared workspace. This is correct behavior: the synthesis contract makes clean workspace available but does not require its use. These small targets qualify the contract and scheduler integration; they are not evidence that the learner has discovered a transformation that is impossible without ancillas.

\subsection{QFT-3 witness certification and search probe}

The independently constructed one-clean-ancilla QFT-3 witness is accepted. Its metrics are
\begin{equation}
\Delta_{\mathrm{proj}}=1.1240\times10^{-15},
\qquad
\ell_A=2.1513\times10^{-32}.
\end{equation}
The residual values are at floating-point roundoff scale. The witness uses 47 native gates, 21 $T/T^\dagger$ gates, and 19 CNOTs, and returns the workspace qubit to $\lvert0\rangle$ for every logical input.

The unrestricted hierarchical search probe does not rediscover a QFT-3 solution within the fixed 1.5-second bound. It stops at the wall limit after 150 outer allocations and 585 exact continuation attempts, with peak frontier 568. Figure~\ref{fig:qft3-probe} shows the recorded work counts.

\begin{figure}[t]
\centering
\includegraphics[width=\columnwidth]{qft3_search_probe.png}
\caption{Bounded unrestricted QFT-3 search probe. The negative result separates representability and contract correctness from practical discovery at the tested search budget.}
\label{fig:qft3-probe}
\end{figure}

Thus, the experiment establishes that the state representation, full-register canonicalizer, isometry propagation, clean-workspace restoration test, and independent certifier can represent a nontrivial ancilla-assisted QFT-3 implementation. It does not establish that the present linear policies can synthesize that 47-gate circuit from unrestricted search within a short wall-clock budget.

\subsection{Component qualifications retained by the ancilla engine}

The strengthened canonicalization audit reports no observed false semantic merges on the bounded sample. Relative to the earlier local key, unresolved duplicate excess decreases by $50.97\%$ for the one-qubit enumeration and $50.0\%$ for the two-qubit enumeration. Median key cost increases by factors $3.81$ and $4.38$, respectively. These measurements justify caching the canonical view and using deferred expansion to avoid constructing it for every sibling.

The mixed Clifford+$T$ continuation suite remains a direct qualification of the two-level scheduler. On its selected four-, five-, and six-qubit targets, frozen LinUCB scheduling uses fewer exact continuations than deferred native ordering and eager outer SARSA. The ancilla extension retains the same policy architecture while adding role-aware contexts and contract-isometry features.
